% Generic, venue-agnostic starting skeleton.
%
% This intentionally uses the plain `article` class rather than a journal's
% own .cls file (REVTeX, JHEP3, svjour3, ...): those are maintained and
% distributed by the publisher/society and should be pulled fresh from the
% journal's site once the target venue is fixed -- see
% references/latex-and-formatting.md for where to get each one and a
% ready-to-adapt skeleton in that venue's actual class.
%
% This file is meant as an early, venue-agnostic draft starting point, or
% for internal notes / theses / arXiv-only preprints that don't need a
% specific journal class.

\documentclass[11pt]{article}

\usepackage[margin=1in]{geometry}
\usepackage{amsmath,amssymb}
\usepackage{graphicx}
\usepackage{booktabs}
\usepackage[colorlinks=true, citecolor=blue, linkcolor=blue, urlcolor=blue]{hyperref}

\title{Title States the Result, Not the Method}
\author{First Author$^{1}$, Second Author$^{2}$ \\[4pt]
  \small $^1$Institution A, City, Country \\
  \small $^2$Institution B, City, Country}
\date{\today}

\begin{document}

\maketitle

\begin{abstract}
  One sentence of context. One to two sentences on the dataset/method used.
  One sentence stating the headline result with its uncertainty. One
  sentence on significance and comparison to prior work or theory. Write
  this section last, after every number below is final.
\end{abstract}

\section{Introduction}
\label{sec:intro}

State the physics motivation, summarize prior measurements with citations
(e.g.~\cite{example:key}), and state explicitly what this paper adds.

\section{Dataset and Method}
\label{sec:dataset}

Describe the data source / theoretical framework in enough detail for an
expert reader to judge adequacy, not to re-derive the setup from scratch.

\section{Analysis}
\label{sec:analysis}

Describe the selection / calculation as a narrative, with a cutflow table
(Table~\ref{tab:cutflow}) if useful.

\begin{table}[htbp]
  \centering
  \caption{Example cutflow table. Replace with real numbers; delete if not
    applicable.}
  \label{tab:cutflow}
  \begin{tabular}{lrr}
    \toprule
    Selection & Events & Efficiency [\%] \\
    \midrule
    Preselection & 000{,}000 & 100.0 \\
    Trigger      & 000{,}000 &  00.0 \\
    Final        & 000{,}000 &  00.0 \\
    \bottomrule
  \end{tabular}
\end{table}

\section{Systematic Uncertainties}
\label{sec:systematics}

Organize by source; summarize in a table (source, size, correlation)
alongside supporting text on how each was evaluated.

\section{Results}
\label{sec:results}

State the headline number plainly, with statistical and systematic
uncertainty distinguished, e.g.\ $125.3 \pm 0.4~\mathrm{(stat)} \pm
0.6~\mathrm{(syst)}$~GeV, and compare explicitly to theory or prior
measurements (Fig.~\ref{fig:result}).

\begin{figure}[htbp]
  \centering
  % \includegraphics[width=0.9\linewidth]{figures/result.pdf}
  \caption{Replace with the actual result figure. Caption should state what
    is shown, briefly how it was made, and the one-sentence takeaway.}
  \label{fig:result}
\end{figure}

\section{Summary}
\label{sec:summary}

Restate the result in one or two sentences, note implications, and
optionally outlook -- no new numbers not already given in
Sec.~\ref{sec:results}.

\section*{Acknowledgments}

Funding agencies and collaboration boilerplate go here; confirm current
wording with the collaboration's publications committee rather than reusing
an old paper's text verbatim.

\bibliographystyle{plain} % swap for the venue's required style, e.g. apsrev4-2 or JHEP
\bibliography{references}

\end{document}
